% Original Girard1987 CUT trips and cyclic-door boxes; earlier packets frozen.
\section{Original CUT trips and cyclic-door proof boxes}
\label{sec:p07-girard-cut-boxes}

We continue the original two-gate trip criterion, switching both tensor
and par. General-formula axiom links retain distinct occurrence identities.
First comes the closed unit-free PN0 cut system, then the separate
opaque-box theorem and actual PN1/PN2 surfaces
\cite[Sections~2.2--2.5, pp.~40--47]{Girard1987AuthorHosted}.

\begin{hbdefinition}{def:p07-girard-cut-marker}{Direct cut and the terminal dummy marker}
An actual cut consumes distinct occurrences \(X,Y\) with dual formula
labels and has no formula conclusion.
Its crossings are \(X^-\mapsto Y^+\), \(Y^-\mapsto X^+\).
Girard's marker representation adds a fresh symbol \(\kappa\), printed
CUT, with two gates. It is treated locally as a tensor root with an
independent \(L/R\) position, and is necessarily terminal.
It is not an object-language formula, cannot be consumed and is not
a logical sequent conclusion. Different cuts have different marker
identities. The exact clauses are on \cref{sheet:p07-girard-cut-ports}.
For \(p\) ordinary occurrences and \(c\) cuts, there are \(2(p+c)\)
marked gates and source size \(s=p+c\).
\end{hbdefinition}

\begin{hblemma}{thm:p07-girard-cut-erasure}
 {proved=annex:p07.girard-cut.erasure}{Dummy gates subdivide cut orbits}
Erasing every marker's gates gives exactly the direct crossings.
At fixed ordinary positions the marked and direct permutations have
corresponding orbits, independently of marker positions.
Original marked correctness is therefore equivalent to one direct
orbit at every ordinary setting.
\end{hblemma}

\begin{hbtheorem}{thm:p07-girard-cut-sequentialization}
 {proved=annex:p07.girard-cut.sequentialization}{Original reconstruction with cut}
A marked PN0 structure is trip-correct exactly when it is the original
link construction of a closed multiplicative sequent proof with cut.
Reconstruction replaces CUTs temporarily by tagged terminal tensors,
then replaces those designated tensor inferences by cuts, using
\cref{thm:p07-girard-empire-sequentialization}.
\end{hbtheorem}

\begin{hbdefinition}{def:p07-girard-cut-contractions}{The original incidence contractions}
For a cut of \(U=B\,m\,C\) and
\(V=B^\perp\,m^\perp\,C^\perp\), delete the two opposite
multiplicative producers, their roots and old CUT marker.
Add cuts on \(B,B^\perp\) and \(C,C^\perp\), with fresh markers.
Both orientations \(m=\otimes\) and \(m=\bindnasrepma\) are allowed.
For an axiom producing cut premise \(X\) and companion \(Z\), with
opposite cut premise \(Y\) carrying \(Z\)'s label, delete that axiom,
its cut-facing occurrence and cut. Glue \(Z,Y\), retaining \(Y\)'s
producer and \(Z\)'s consumer or terminal boundary.
The mirrored case is allowed, and either choice may apply when both
cut premises have axiom producers. This is occurrence gluing, not
identification of every equal formula label
\cite[pp.~41--42]{Girard1987AuthorHosted}.
\end{hbdefinition}

\begin{hbtheorem}{thm:p07-girard-cut-principal-preservation}
 {proved=annex:p07.girard-cut.principal-preservation}{Both principal orientations preserve original trips}
Every principal contraction of a trip-correct PN0 structure is
trip-correct. For an arbitrary reduct setting the proof partitions the
old trip into four disjoint outside slices and rejoins all four in one
orbit. Arbitrary new marker positions are included.
\end{hbtheorem}

\begin{hbtheorem}{thm:p07-girard-cut-axiom-preservation}
 {proved=annex:p07.girard-cut.axiom-preservation}{All axiom-cut cases preserve original trips}
An axiom contraction, its mirrored case and either choice at a
two-axiom cut preserve original correctness. The same-axiom loop is
excluded by its explicit short trips, not by assuming that the
companion occurrences were distinct.
\end{hbtheorem}

\begin{hbtheorem}{thm:p07-girard-cut-normalization}
 {proved=annex:p07.girard-cut.normalization}{Original size decrease and PN0 normalization}
A principal step decreases \(s=p+c\) by one; an axiom step by three.
Every cut in a correct PN0 structure admits one of these steps.
Every maximal reduction ends cut-free in fewer than \(s\) steps.
This occurrence-plus-CUT measure is not a box-duplication bound.
\end{hbtheorem}

\begin{hbtheorem}{thm:p07-girard-cut-confluence}
 {proved=annex:p07.girard-cut.confluence}{All original contraction overlaps}
PN0 contraction is confluent up to incidence isomorphism preserving
retained conclusion identities, labels and connective port orders.
The proof covers disjoint cuts, two cuts sharing an axiom, and both
axiom choices at one cut. It yields the original unique cut-free
normal form, not a unique sequent proof tree
\cite[Proposition~2.11 and Theorem~2.12, p.~42]{Girard1987AuthorHosted}.
\end{hbtheorem}

\begin{hbexample}{def:p07-girard-cut-loop}{The same-axiom short trips}
Cut the two outputs \(X,Y=X^\perp\) of one axiom together.
The direct permutation has two orbits,
\(X^+,Y^-,X^+\) and \(Y^+,X^-,Y^+\).
A marker subdivides one orbit, leaving the other short.
This is not a proof-net and cannot invalidate the strict size theorem.
\end{hbexample}

\begin{hbexample}{def:p07-girard-cut-slices}{An exhaustive principal trip computation}
Cut \(U=B\otimes C\) against \(V=B^\perp\bindnasrepma C^\perp\).
With both producer positions \(L\), let \(B',C'\) name the distinct
dual premises. The four outside slices are
\(P_C:C^+\ldots C^-\), \(P_B:B^+\ldots B^-\),
\(Q_1:B'^+\ldots C'^-\), \(Q_2:C'^+\ldots B'^-\).
The old direct orbit orders them \(Q_1,Q_2,P_C,P_B\), with the root
gates on the intervening crossings; the reduct orders them
\(Q_1,P_C,Q_2,P_B\). The proof establishes that these disjoint slices
cover every remaining gate, not merely that a drawing can be rearranged.
\end{hbexample}

\section{Opaque boxes and all cyclic door positions}
\label{sec:p07-girard-opaque-boxes}

\begin{hbdefinition}{def:p07-girard-box-doors}{The original black-box interface}
A box \(\mathcal B\) produces \(n\geq1\) distinct formula doors
\(F_1,\ldots,F_n\), with no incoming occurrence in the exterior graph.
Its position is any one-cycle permutation \(\rho\) of the door
indices, not an arbitrary permutation and not just rotations of a
preferred order. It routes \(F_i^+\mapsto F_{\rho(i)}^-\).
For \(n=1\) the only position is the identity.
The exterior must have one long orbit at every ordinary setting and
every allowed cyclic position of every box. The actual internal
formation and correctness are checked separately
\cite[Definition~2.14, pp.~43--44]{Girard1987AuthorHosted}.
\end{hbdefinition}

\begin{hblemma}{thm:p07-girard-box-slices}
 {proved=annex:p07.girard-box.slices}{Every exterior slice returns to its own door}
Fix the positions outside one box in a correct exterior.
Deleting its \(n\) successor arcs leaves exactly the paths
\(F_i^-\ldots F_i^+\), covering all gates.
An outside path to a different door would form a proper short trip
at another allowed cyclic position.
\end{hblemma}

\begin{hblemma}{thm:p07-girard-box-empire}
 {proved=annex:p07.girard-box.empire}{The full box analogue of axiom closure}
For a tensor-premise empire in a correct boxed exterior,
membership of one door implies membership of all doors of that box.
The proof includes a distinguished premise that is itself a door,
both same-slice orientations, different-slice endpoints and \(n=1\).
\end{hblemma}

\begin{hbtheorem}{thm:p07-girard-box-sequentialization}
 {proved=annex:p07.girard-box.sequentialization}{Original Theorem 2.15 for opaque exteriors}
Correct multiplicative exteriors with boxes are exactly the link
constructions of sequent proofs whose additional named axiom leaves
are the boxes' entire door sequents. Each box occurrence is used once
as its designated leaf, never split or copied across a tensor inference.
The cut-free result extends to CUT markers by the tagged-tensor
construction. An unchecked leaf is not thereby an admitted logical axiom.
\end{hbtheorem}

\begin{hbtheorem}{thm:p07-girard-box-plugging}
 {proved=annex:p07.girard-box.plugging}{The precise same-door replacement guarantee}
Filling a correct exterior box with a correct original net having
exactly the same nonempty terminal door list, labels and identities
gives a correct plugged net. Internal CUT markers are not extra doors.
The inner trip induces an allowed cyclic permutation, and exterior
jumps expand into disjoint inner slices. This is a one-net filling;
a with-box's pair of certificates has a different formation regime.
\end{hbtheorem}

\begin{hbexample}{def:p07-girard-box-cycles}{Cyclic and one-door boundaries}
A bare terminal three-door box has one six-gate orbit at either
three-cycle position, alternating box jumps with terminal returns.
Allowing the identity would give three two-gate orbits.
Thus \emph{all permutations} would be the wrong criterion.
A one-door box has the orbit \(F^+,F^-,F^+\), the source unit-axiom
interface rather than a bare nullary vertex in a binary unit-free graph.
\end{hbexample}

\begin{hbexample}{def:p07-girard-box-normalization-boundary}{Correct exteriors may have no PN0 contraction}
Two opaque one-door boxes labeled \(a,a^\perp\), joined by a cut,
have one direct four-gate orbit; the marker only subdivides it.
Neither premise has an ordinary axiom or multiplicative producer,
so no PN0 contraction applies. These leaves are formal additional
axioms, not established one-formula proofs of \(a,a^\perp\).
Exterior correctness, interior validation and box-specific normalization
are separate claims.
\end{hbexample}

\section{The actual PN1 and PN2 formation surfaces}
\label{sec:p07-girard-pn1-pn2}

\begin{hbdefinition}{def:p07-girard-box-front-aux}{Front doors, auxiliary contexts and certificates}
PN1's logical-axiom box has front doors \(A,A^\perp\);
the unit axiom has front \(1\); the top axiom has front \(\top\)
and auxiliary context \(\Gamma\).
The finite occurrence list \(\Gamma\) represents the source's other
conclusions; CUT symbols are excluded from door lists.
Write \(\operatorname{Conc}(\beta)\) for that ordinary terminal list.
A with-box has front \(A\mathbin{\&}B\), auxiliary \(\Gamma\),
and separately correct inner nets on \(\Gamma,A\) and \(\Gamma,B\).
Their corresponding auxiliary occurrences match one outer context by
the formation rule, not by raw disjoint tensor gluing.
Bottom and quest-weakening boxes retain an actual inner net on
\(\Gamma\) and add front \(\bot\) or \(?A\).
Promotion requires inner conclusions \(A,\Gamma_?\), every auxiliary
formula quest-prefixed, and has front \(!A\).
The exact cards are on \cref{sheet:p07-girard-pn1-pn2}
\cite[pp.~44--46]{Girard1987AuthorHosted}.
\end{hbdefinition}

\begin{hbdefinition}{def:p07-girard-box-unary}{Unary links and quest contraction}
The unary links are both plus injections and dereliction.
A link with premise \(F\), root \(G\) routes
\(G^+\mapsto F^+\), \(F^-\mapsto G^-\), without a switch.
Quest contraction consumes two distinct occurrences of \(?A\) and
produces a third \(?A\) occurrence, with the geometry and both
positions of a par. Equal labels do not identify those occurrences.
This is not unrestricted structural contraction.
\end{hbdefinition}

\begin{hblemma}{thm:p07-girard-unary-insertion}
 {proved=annex:p07.girard-box.unary-insertion}{Local original-trip unary introduction}
Appending a source unary link to a terminal premise preserves one
long trip at every position, inserting its two root gates into that
terminal return. Removing a terminal unary root reverses it.
This is not the full PN1/PN2 normalization theorem.
\end{hblemma}

\begin{hbdefinition}{def:p07-girard-pn2-substitution}{The second-order substitution and source rules}
PN2 uses propositional \(\forall a,\exists a\), the source's large
wedge and vee, not first-order binders.
Capture-avoiding \(A[B/a]\) replaces free \(a\) by \(B\) and free
\(a^\perp\) by \(B^\perp\), first renaming binders as necessary.
The existential link is unary from \(A[B/a]\) to \(\exists a.A\).
The universal box requires an actual inner net on \(\Gamma,A\),
with \(a\) not free in \(\Gamma\), and has front \(\forall a.A\)
and auxiliary \(\Gamma\). Freshness is part of the rule, not a
consequence of a generic cyclic box
\cite[pp.~46--47]{Girard1987AuthorHosted}.
\end{hbdefinition}

\begin{proofboundary}{def:p07-girard-box-full-source-boundary}{Formation versus whole normalization}
The registry retains every PN1/PN2 surface of Sections 2.4--2.5.
The opaque-box theorem proves the exterior representation with named
leaves; those cards still require their actual inner certificates and
side conditions. It does not prove every box duplication, quantifier,
symmetric or commutative contraction and its normalization argument
in Section 4. In particular a with-box's pair is not a one-net filling,
and a cut against arbitrary doors is not a PN0 axiom contraction.
The source's unboxed-weakening, quantifier-box removal and additive-box
identification discussions retain their limited or unresolved standing
\cite[p.~47]{Girard1987AuthorHosted}. The exact 1989 presentation and
GoI remain separate requirements.
\end{proofboundary}

\showcalculus{p07-girard-cut-boxes}
\index{proof net!original CUT trips}
\index{proof box!cyclic door permutation}
